\section{Actual native model potentials and finite rational FS minorants}

Use the actual original test configuration, original total invertible
line $M$, fibre-one line $L$, action and fibre-one identifications.
All values and norms refer to these original objects.

\begin{definition}\label{def:native-tc-original-model-potential}
\lean{EnvelopeProperty.NativeTestConfiguration.TestConfiguration.gaussNativeModelPotential,EnvelopeProperty.NativeTestConfiguration.TestConfiguration.gaussNativeModelPotential_eq_compatibleFactor}
\leanok
For every $0<r\leq1$ and actual original native point $x$, define
\[
 \phi_{\mathcal X,L}(x,r)=-\log c_{\mathcal X,L}(x,r),
\]
where the factor is the actual native valuation of the original
geometric frame-transition unit. It is a finite real-valued function.
For ample $L$ its values are independent of all original center-open
frame choices.
\uses{def:native-tc-intrinsic-native-unit-factor,thm:native-tc-intrinsic-factor-original-frame-independence}
\end{definition}
\begin{proof}
The original geometric frames define a genuine unit in the actual
product residue field. Its native valuation is strictly positive,
so its real logarithm is finite. Original ampleness proves the
factor equals the factor formed from every other original
center-open frame. Applying the same real logarithm and negation
proves the asserted independence of this function.
\end{proof}

\begin{theorem}\label{thm:native-tc-original-potential-weighted-fs}
\lean{EnvelopeProperty.NativeTestConfiguration.TestConfiguration.gaussPowerFamilyNorm_eq_nativeFactor,EnvelopeProperty.NativeTestConfiguration.TestConfiguration.gaussNativeModelPotential_weightedFSData_eval}
\leanok
For every positive original degree $n$, every actual original finite
family $G$, every $0<r\leq1$ and every original point $x$,
\[
 \operatorname{FS}_{n,G,r}(x)=\phi_{\mathcal X,L}(x,r)
 +\frac1n\log\max_{s\in G}\lVert s\rVert_{\mathcal X,r,x}.
\]
The same actual native factor also computes the whole finite-family
norm in all degrees, including zero, as $c^n$ times the maximum of
the actual weighted Laurent coefficient norms.
\uses{def:native-tc-original-model-potential,cor:native-tc-coherent-all-power-family-norms,thm:native-tc-original-weighted-fs-evaluation}
\end{theorem}
\begin{proof}
Take the finite maximum in the actual all-degree all-section norm
formula for the defined native factor. Multiplication by $c^n$
preserves finite maxima, and taking the union of original Laurent
supports is exactly the actual finite coefficient/exponent family.
Taking the extended logarithm and dividing by positive $n$ gives
$\log c+\operatorname{FS}_{n,G,r}$. Add the finite real $-\log c$
to this equality. Cancellation gives the displayed identity, also
when a family is empty or vanishes at the point.
\end{proof}

\begin{theorem}\label{thm:native-tc-original-potential-finite-rational-fs-minorants}
\lean{EnvelopeProperty.NativeTestConfiguration.TestConfiguration.exists_eventual_finiteFS_modelPotential_minorants}
\leanok
Suppose the original fibre-one line is ample and put $r=e^{-1}$.
There is an original integer $N$ such that for every positive
$n\geq N$ an actual original finite family $G_n$, chosen before
all points, defines a genuine finite rational FS function of the
original $L$ and original class one, and
\[
 \operatorname{FS}_{n,G_n,r}(x)
 =\phi_{\mathcal X,L}(x,r)+\frac1n
   \log\lVert\mathfrak a_n\rVert_{\mathcal X,r,x}
 \leq\phi_{\mathcal X,L}(x,r)
\]
for every actual original point. Here $\mathfrak a_n$ is the base
ideal of all original global sections of $M^{\otimes n}$.
\uses{thm:native-tc-original-potential-weighted-fs,thm:native-test-configuration-actual-gauss-finite-norm-family,cor:native-tc-original-rational-fs-families,def:native-tc-original-normalized-log-ideal}
\end{theorem}
\begin{proof}
Actual original ampleness supplies the finite family in each eventual
degree, simultaneously computing the entire base-ideal norm at all
points with strictly positive finite maximum. Substitute this norm
identity in the preceding formula. Its logarithm is finite, so the
FS value is not negative infinity. The actual native FS nonvanishing
criterion then supplies a nonzero section norm at every point,
proving that there is no common native zero. The original Laurent
weights are integers times $\log(e^{-1})=-1$, so they are rational;
these actual data give the stated finite rational FS function.
Finally every actual model ideal norm is at most one. Its extended
logarithm is nonpositive, and division by the positive original
degree preserves this inequality. Adding the finite actual model
potential proves the displayed bound.
\end{proof}

\begin{remark}
These are genuine original FS minorants of the defined native model
potential. The source Cartier-model identification and the reverse
section-extension comparison for the final envelope formula remain
separate geometric steps.
\end{remark}
